\documentclass[12pt]{article}
\usepackage[T1]{fontenc}
\usepackage[utf8]{inputenc}
\usepackage{lmodern}
\usepackage{textcomp}
\usepackage{enumitem}
\usepackage{geometry}
\geometry{margin=1in}
\begin{document}
\begin{center}
Answer Key - Science - K-12 Level Assessment 20 \
\small (Answers are ordered exactly as in the question paper.)
\end{center}

\section*{Multiple Choice}
\begin{enumerate}[label=\arabic*.]
\item \textbf{Answer:} C
\\ \textit{Options:} A) visible color; B) composition; C) surface texture; D) temperature
\\ \textit{Note:} Surface texture is the least useful criterion for classifying stars because it is not a fundamental property that significantly affects a star's characteristics or classification compared to factors like luminosity, temperature, or spectral type.
\item \textbf{Answer:} A
\\ \textit{Options:} A) Earth is heated by the Sun.; B) Mountains direct moist air upward.; C) Water forms droplets when cooled.; D) Most of Earth is covered with water.
\\ \textit{Note:} The primary cause of rainstorms is that the Sun heats the Earth's surface, causing water to evaporate and rise into the atmosphere, where it cools and condenses to form clouds that eventually release precipitation.
\item \textbf{Answer:} B
\\ \textit{Options:} A) a hill and a delta; B) a delta and a canyon; C) a canyon and a dune; D) a dune and a delta
\\ \textit{Note:} The river creates a delta at its mouth where it deposits sediment, while the continuous erosion of its banks over time can lead to the formation of a canyon along its course.
\item \textbf{Answer:} C
\\ \textit{Options:} A) research; B) interviews; C) observation; D) measurement
\\ \textit{Note:} The correct answer is observation because the student will need to visually assess and record the different colors of light that emerge from the prism after refraction occurs.
\item \textbf{Answer:} A
\\ \textit{Options:} A) They are solids at room temperature.; B) They don't conduct electricity.; C) They are brittle and dull.; D) They are radioactive.
\\ \textit{Note:} The correct answer is based on the fact that most elements on the left side of the Periodic Table, particularly the metals, are indeed solid at room temperature, with the exception of mercury, which is a liquid.
\end{enumerate}

\section*{True / False}
\begin{enumerate}[label=\arabic*.]
\setcounter{enumi}{5}
\item \textbf{Answer:} False
\\ \textit{Options:} A) True; B) False
\\ \textit{Note:} Scientific discoveries undergo rigorous peer review and critical evaluation by the scientific community before being accepted as valid.
\item \textbf{Answer:} False
\\ \textit{Options:} A) True; B) False
\\ \textit{Note:} Reusing old materials often leads to lower costs for households because it reduces the need to purchase new items, minimizes waste, and can lead to savings in disposal fees.
\item \textbf{Answer:} False
\\ \textit{Options:} A) True; B) False
\\ \textit{Note:} Early photosynthetic life, such as cyanobacteria, actually contributed to the rise of oxygen in the atmosphere and did not significantly decrease carbon dioxide levels until much later in Earth's history.
\item \textbf{Answer:} False
\\ \textit{Options:} A) True; B) False
\\ \textit{Note:} While water's ability to freeze is important, its other properties, such as being a solvent, regulating temperature, and supporting cellular functions, are more crucial for sustaining life on Earth.
\item \textbf{Answer:} True
\\ \textit{Options:} A) True; B) False
\\ \textit{Note:} The kidneys filter blood to eliminate waste products and excess substances, producing urine that carries these dissolved waste molecules out of the body.
\end{enumerate}

\section*{Long Form Response}
\begin{enumerate}[label=\arabic*.]
\setcounter{enumi}{10}
\item \textbf{Answer:} A culture living above the Arctic Circle might develop inventions such as insulated clothing, effective heating systems, and specialized transportation like sleds or snowmobiles to adapt to their cold environment. The extreme weather conditions, including long winters and limited daylight, would influence these innovations, leading to a focus on warmth, mobility, and efficient use of resources. Additionally, tools for hunting and fishing, such as ice augers and snowshoes, would be essential for obtaining food in harsh conditions. These inventions demonstrate how the Arctic environment shapes the needs and creativity of its inhabitants, allowing them to thrive despite the challenges they face.
\\ \textit{Note:} correct because it identifies specific inventions that address the challenges of extreme cold and limited resources in the Arctic environment, emphasizing the need for warmth, mobility, and tools for subsistence activities like hunting and fishing.
\item \textbf{Answer:} During a nature walk, students could make several key observations about grasshoppers, such as their physical characteristics, behaviors, and interactions with their environment. For instance, they might notice the grasshopper's long hind legs, which are adapted for jumping, and how these adaptations allow them to escape predators. Additionally, observing their feeding habits on grasses can help students understand their role in the ecosystem as herbivores. Analyzing these behaviors and traits is significant because it helps students appreciate how grasshoppers contribute to the food web and the balance of their habitats, fostering a deeper understanding of ecological relationships.
\\ \textit{Note:} correct because it identifies critical observations of grasshoppers' physical adaptations and behaviors, illustrating their ecological role and enhancing students' understanding of predator-prey dynamics and the importance of herbivores in their environment.
\end{enumerate}

\vfill
\noindent\textit{End of Answer Key}
\end{document}